\documentclass[11pt]{article}
\usepackage[margin=1in]{geometry}
\usepackage{amsmath,amssymb}
\usepackage{booktabs}
\usepackage{enumitem}
\usepackage{hyperref}
\usepackage{graphicx}

\title{Stage~1 Workflow Summary:\\
  From Raw Data to Reference-Site Threshold Selection}
\author{}
\date{}

\begin{document}
\maketitle

%% ----------------------------------------------------------------
\section{Overview}

This document summarises the three-step workflow that transforms raw
sediment-chemistry, environmental, and benthic-invertebrate data into
(i)~a site-level contamination score, (ii)~an ecologically validated
scoring rule, and (iii)~a defensible reference-site threshold.

%% ----------------------------------------------------------------
\section{Step 1: PCA-Based Contamination Scoring}

\begin{enumerate}[leftmargin=*]
  \item \textbf{Input.}
    A 233-site $\times$ 56-variable workbook containing 16~chemical
    contaminants, 5~environmental variables, and 16~benthic taxa
    (octave-transformed abundances).

  \item \textbf{Log-transformation.}
    Chemical concentrations are transformed via $\log_2(1+x)$ to
    reduce right skew.

  \item \textbf{PCA.}
    Principal Component Analysis is fitted on the 16 log-transformed
    chemical variables with min--max standardisation.  The first five
    PCs are retained (cumulative variance $\geq 83\%$).

  \item \textbf{Sign orientation.}
    Each PC is oriented so that higher scores correspond to higher
    contamination (positive correlation with known contaminant
    indicators).

  \item \textbf{Component rescaling.}
    Each retained PC score is min--max rescaled to $[0,1]$.

  \item \textbf{Two scoring rules.}
    \begin{itemize}
      \item \textbf{SumRel} — sum of rescaled PC scores.
        Reflects cumulative, multi-dimensional contamination pressure.
      \item \textbf{MaxRel} — maximum of rescaled PC scores.
        Captures the single worst contamination axis per site.
    \end{itemize}
\end{enumerate}

\noindent
\textbf{Outputs:}
site-level SumRel and MaxRel scores, PC loadings, variance-explained
plot, ridge-loading plot, and corridor-bifurcation maps for each score.

%% ----------------------------------------------------------------
\section{Step 2: Ecological Validation via RDA}

The competing scores are evaluated by how well the resulting
``least-contaminated'' reference subset explains taxa composition
through environmental gradients.

\begin{enumerate}[leftmargin=*]
  \item \textbf{Reference-site selection.}
    For a given threshold proportion $q$ (e.g.\ 25\%), the bottom-$q$
    fraction of sites ranked by the contamination score is designated
    as the reference subset.

  \item \textbf{Redundancy Analysis (RDA).}
    Within the reference subset, taxa (response, $\mathbf{Y}$) are
    regressed on z-scored environmental variables (predictors,
    $\mathbf{X}$).  RDA decomposes explained variation into
    constrained axes (RDA1, RDA2, \ldots).

  \item \textbf{Permutation tests.}
    A global pseudo-$F$ test (999~permutations) assesses overall
    model significance.  Per-axis and per-term tests identify which
    axes and environmental variables contribute significantly.

  \item \textbf{Score comparison.}
    At each tested threshold, SumRel consistently yields higher
    adjusted~$R^2$ and lower $p$-values than MaxRel, indicating that
    the cumulative scoring rule selects reference sites where
    environmental gradients more cleanly explain biological variation.
\end{enumerate}

%% ----------------------------------------------------------------
\section{Step 3: Reference-Threshold Sensitivity Analysis}

Rather than fix the threshold at an arbitrary value, a grid sweep
from 5\% to 99\% (step~2\%) evaluates how RDA performance changes
with reference-subset size.

\begin{enumerate}[leftmargin=*]
  \item \textbf{Metrics collected at each threshold:}
    $R^2$, adjusted~$R^2$, global pseudo-$F$, $p$-value, constrained
    and total inertia, RDA1 eigenvalue, and sample size.

  \item \textbf{Identifying the recommended range.}
    Adjusted~$R^2$ for SumRel shows a local plateau between
    approximately 22\%--28\%, before a steep decline at larger
    thresholds where contaminated sites dilute the reference pool.
    Within this range the global test remains highly significant
    ($p = 0.001$) and sample sizes exceed 50, ensuring statistical
    power.

  \item \textbf{$F$-critical threshold.}
    On the pseudo-$F$ panel, the $F$-statistic corresponding to
    $p = 0.05$ and the threshold at which it is first reached are
    identified by linear interpolation and reported in the legend.

  \item \textbf{Selected threshold: 25\%.}
    A 25\% cutoff sits in the middle of the recommended plateau,
    yielding $\sim$58 reference sites with adjusted~$R^2 \approx 0.22$
    and $p = 0.001$.
\end{enumerate}

%% ----------------------------------------------------------------
\section{Final Decisions}

\begin{itemize}
  \item \textbf{Scoring rule:} SumRel (ecological view --- cumulative
    contamination pressure).
  \item \textbf{Reference threshold:} 25\% of least-contaminated
    sites.
  \item \textbf{Full RDA} is fitted at the chosen threshold to produce
    the definitive triplot, axes-summary, and terms-summary tables.
\end{itemize}

%% ----------------------------------------------------------------
\section{Code Organisation}

All computation is implemented in the \texttt{src/zci} package:

\begin{description}[style=nextline]
  \item[\texttt{run\_stage1.py}]
    Runs PCA + dual scoring (SumRel / MaxRel).
  \item[\texttt{run\_stage\_threshold.py}]
    Runs Stage~1, sweeps thresholds, fits full RDA at 25\%, and saves
    sensitivity plots, comparison tables, triplot, and RDA summary
    tables.
\end{description}

\noindent
The standalone \texttt{run\_stage\_rda.py} is superseded by the
integrated threshold-sensitivity pipeline, which performs the same
RDA fitting at the selected threshold while also providing the full
sensitivity analysis.

\end{document}
